% tonghopquizz.tex — Tự động chuyển đổi từ tonghopquizz.md (157 câu).
% Biên dịch: xelatex tonghopquizz.tex (chạy 2 lần để cập nhật mục lục).
% Sử dụng XeLaTeX để hiển thị tiếng Việt UTF-8 và các ký hiệu kỹ thuật.
\documentclass[11pt,a4paper]{article}
\usepackage{fontspec}
\IfFontExistsTF{Noto Serif}{\setmainfont{Noto Serif}}{
  \IfFontExistsTF{DejaVu Serif}{\setmainfont{DejaVu Serif}}{\setmainfont{Latin Modern Roman}}
}
\usepackage[vietnamese]{babel}
\usepackage[a4paper,top=21mm,bottom=22mm,left=23mm,right=22mm,headheight=16pt]{geometry}
\usepackage{xcolor}
\usepackage{needspace}
\usepackage{fancyhdr}
\usepackage{xurl}
\usepackage{hyperref}
\usepackage{microtype}
\hypersetup{colorlinks=true,linkcolor=black,urlcolor=blue,pdfauthor={Tong hop cau hoi IoT},pdftitle={Tong hop 157 cau hoi IoT}}
\definecolor{answer}{HTML}{123C69}
\definecolor{sourcegray}{HTML}{5C6570}
\setlength{\parindent}{0pt}
\setlength{\parskip}{3pt}
\setlength{\emergencystretch}{3em}
\raggedbottom
\pagestyle{fancy}
\fancyhf{}
\fancyhead[L]{\small Cơ sở và thiết kế hệ thống IoT}
\fancyhead[R]{\small Ngân hàng câu hỏi}
\fancyfoot[C]{\thepage}
\renewcommand{\headrulewidth}{0.3pt}
\begin{document}
\begin{center}
{\LARGE\bfseries TỔNG HỢP CÂU HỎI ÔN TẬP IoT}\\[8pt]
{\large 157 câu hỏi và đáp án đã rà soát}\\[4pt]
{\small Cập nhật ngày 14/10/2026 — Tài liệu xuất từ \texttt{tonghopquizz.md}}
\end{center}

\noindent\textit{Hướng dẫn:} Mỗi câu hỏi được trình bày kèm đáp án \textbf{in đậm}.
Các phương án nhiễu đã bị loại bỏ trong tài liệu Markdown nguồn sẽ không được tự tạo lại.
Những đáp án do AI hiệu đính hoặc còn cần đối chiếu đều có ghi chú riêng.

\tableofcontents
\bigskip

\section*{I. Công nghiệp 4.0 và chuyển đổi số}
\addcontentsline{toc}{section}{I. Công nghiệp 4.0 và chuyển đổi số}

\needspace{5\baselineskip}
\textbf{Câu 1.} Cách mạng công nghiệp 4.0 đề cập đến sự kết hợp giữa công nghệ nào?\par
{\color{answer}\textbf{Đáp án: Công nghệ vật lý, công nghệ số và sinh học.}}\par
{\footnotesize\color{sourcegray}\textit{Nguồn:} GK100 — câu 1.\par}
\vspace{3pt}\par\noindent{\color{black!18}\rule{\linewidth}{0.3pt}}\par\medskip

\needspace{5\baselineskip}
\textbf{Câu 2.} Cuộc cách mạng công nghiệp 4.0 được định nghĩa bởi ai?\par
{\color{answer}\textbf{Đáp án: Klaus Schwab.}}\par
{\footnotesize\color{sourcegray}\textit{Nguồn:} GK100 — câu 2.\par}
\vspace{3pt}\par\noindent{\color{black!18}\rule{\linewidth}{0.3pt}}\par\medskip

\needspace{5\baselineskip}
\textbf{Câu 3.} Công nghệ nào được xem là trụ cột trong cách mạng công nghiệp 4.0?\par
{\color{answer}\textbf{Đáp án: Trí tuệ nhân tạo (AI).}}\par
{\footnotesize\color{sourcegray}\textit{Nguồn:} GK100 — câu 3.\par}
\vspace{3pt}\par\noindent{\color{black!18}\rule{\linewidth}{0.3pt}}\par\medskip

\needspace{5\baselineskip}
\textbf{Câu 4.} Mục tiêu chính của cách mạng công nghiệp 4.0 là gì?\par
{\color{answer}\textbf{Đáp án: Tăng cường sự kết nối và tương tác giữa các yếu tố trong sản xuất.}}\par
{\footnotesize\color{sourcegray}\textit{Nguồn:} GK100 — câu 4.\par}
\vspace{3pt}\par\noindent{\color{black!18}\rule{\linewidth}{0.3pt}}\par\medskip

\needspace{5\baselineskip}
\textbf{Câu 5.} Lựa chọn không phản ánh đặc trưng của CMCN 3.0\par
{\color{answer}\textbf{Đáp án: Internet kết nối vạn vật.}}\par
{\footnotesize\color{sourcegray}\textit{Nguồn:} GK100 — câu 5.\par}
\vspace{3pt}\par\noindent{\color{black!18}\rule{\linewidth}{0.3pt}}\par\medskip

\needspace{5\baselineskip}
\textbf{Câu 6.} Trong bối cảnh của cuộc cách mạng công nghiệp lần thứ 4, khái niệm nào mô tả chính xác nhất sự thay đổi công nghệ?\par
{\color{answer}\textbf{Đáp án: Là cuộc cách mạng sử dụng công nghệ số như IoT, AI, Big Data kết hợp với các ngành công nghiệp khác..}}\par
{\footnotesize\color{sourcegray}\textit{Nguồn:} GK100 — câu 6.\par}
\vspace{3pt}\par\noindent{\color{black!18}\rule{\linewidth}{0.3pt}}\par\medskip

\needspace{5\baselineskip}
\textbf{Câu 7.} Cuộc cách mạng công nghiệp lần thứ nhất là cuộc cách mạng nào?\par
{\color{answer}\textbf{Đáp án: Là cuộc cách mạng cơ khí hóa do động cơ hơi nước gây ra..}}\par
{\footnotesize\color{sourcegray}\textit{Nguồn:} GK100 — câu 7.\par}
\vspace{3pt}\par\noindent{\color{black!18}\rule{\linewidth}{0.3pt}}\par\medskip

\needspace{5\baselineskip}
\textbf{Câu 8.} Cuộc cách mạng công nghiệp lần thứ 4 (CMCN 4.0) có tác động như thế nào đến mô hình sản xuất và kinh doanh?\par
{\color{answer}\textbf{Đáp án: Thay đổi mô hình sản xuất và kinh doanh truyền thống sang mô hình hiện đại, thay đổi thói quen, suy nghĩ của người tiêu dùng..}}\par
{\footnotesize\color{sourcegray}\textit{Nguồn:} GK100 — câu 8.\par}
\vspace{3pt}\par\noindent{\color{black!18}\rule{\linewidth}{0.3pt}}\par\medskip

\needspace{5\baselineskip}
\textbf{Câu 9.} Kỷ nguyên chuyển đổi kỹ thuật số (DX) có tác động như thế nào đến các doanh nghiệp?\par
{\color{answer}\textbf{Đáp án: Chuyển đổi toàn diện các quy trình kinh doanh, mô hình kinh doanh và cách thức tương tác với khách hàng thông qua việc áp dụng công nghệ số..}}\par
{\footnotesize\color{sourcegray}\textit{Nguồn:} GK100 — câu 9.\par}
\vspace{3pt}\par\noindent{\color{black!18}\rule{\linewidth}{0.3pt}}\par\medskip

\needspace{5\baselineskip}
\textbf{Câu 10.} Chuyển đổi số là gì?\par
{\color{answer}\textbf{Đáp án: Là quá trình sử dụng công nghệ số để thay đổi toàn diện cách thức hoạt động, mô hình kinh doanh và cách thức tương tác với khách hàng..}}\par
{\footnotesize\color{sourcegray}\textit{Nguồn:} GK100 — câu 10.\par}
\vspace{3pt}\par\noindent{\color{black!18}\rule{\linewidth}{0.3pt}}\par\medskip

\needspace{5\baselineskip}
\textbf{Câu 11.} Doanh nghiệp số là gì?\par
{\color{answer}\textbf{Đáp án: Là doanh nghiệp sử dụng công nghệ số để tăng cường hiệu quả kinh doanh và tạo ra giá trị mới..}}\par
{\footnotesize\color{sourcegray}\textit{Nguồn:} GK100 — câu 11.\par}
\vspace{3pt}\par\noindent{\color{black!18}\rule{\linewidth}{0.3pt}}\par\medskip

\needspace{5\baselineskip}
\textbf{Câu 12.} Công nghệ M2M (Machine-to-Machine) là gì?\par
{\color{answer}\textbf{Đáp án: Công nghệ kết nối các thiết bị máy móc với nhau qua mạng lưới điện thoại di động hoặc internet mà không cần sự can thiệp của con người.}}\par
{\footnotesize\color{sourcegray}\textit{Nguồn:} GK100 — câu 12.\par}
\vspace{3pt}\par\noindent{\color{black!18}\rule{\linewidth}{0.3pt}}\par\medskip

\needspace{5\baselineskip}
\textbf{Câu 13.} Cách mạng công nghiệp 4.0 đã có một trong những tác động nào đến ngành công nghiệp ô tô?\par
{\color{answer}\textbf{Đáp án: Sản xuất động cơ bằng máy in 3D.}}\par
{\footnotesize\color{sourcegray}\textit{Nguồn:} GK100 — câu 13.\par}
\vspace{3pt}\par\noindent{\color{black!18}\rule{\linewidth}{0.3pt}}\par\medskip

\needspace{5\baselineskip}
\textbf{Câu 14.} Những công nghệ nào dùng để phân tích dữ liệu trong cuộc cách mạng công nghiệp 4.0?\par
{\color{answer}\textbf{Đáp án: Cả ba đáp án.}}\par
{\footnotesize\color{sourcegray}\textit{Nguồn:} GK100 — câu 14.\par}
\vspace{3pt}\par\noindent{\color{black!18}\rule{\linewidth}{0.3pt}}\par\medskip

\needspace{5\baselineskip}
\textbf{Câu 15.} Ứng dụng chính của IoT không bao gồm:\par
{\color{answer}\textbf{Đáp án: Đào tạo trí tuệ nhân tạo.}}\par
{\footnotesize\color{sourcegray}\textit{Nguồn:} GK100 — câu 15.\par}
\vspace{3pt}\par\noindent{\color{black!18}\rule{\linewidth}{0.3pt}}\par\medskip

\needspace{5\baselineskip}
\textbf{Câu 16.} IoT có thể áp dụng trong lĩnh vực nào sau đây?\par
{\color{answer}\textbf{Đáp án: Tất cả các phương án đều đúng.}}\par
{\footnotesize\color{sourcegray}\textit{Nguồn:} GK100 — câu 16.\par}
\vspace{3pt}\par\noindent{\color{black!18}\rule{\linewidth}{0.3pt}}\par\medskip

\needspace{5\baselineskip}
\textbf{Câu 17.} Trong cụm từ "ICBAMS", ký tự 'C' đại diện cho khái niệm gì?\par
{\color{answer}\textbf{Đáp án: Computer (Máy tính), theo đáp án in trong PDF.}}\par
{\footnotesize\color{sourcegray}\textit{Nguồn:} GK100 — câu 17.\par}
{\footnotesize\color{sourcegray}\textit{Ghi chú AI:} Chưa tìm được định nghĩa chuẩn của viết tắt ICBAMS trong học liệu cung cấp; cần giảng viên xác nhận.\par}
\vspace{3pt}\par\noindent{\color{black!18}\rule{\linewidth}{0.3pt}}\par\medskip

\section*{II. Phần cứng, Raspberry Pi và linh kiện điện tử}
\addcontentsline{toc}{section}{II. Phần cứng, Raspberry Pi và linh kiện điện tử}

\needspace{5\baselineskip}
\textbf{Câu 18.} Raspberry Pi là gì?\par
{\color{answer}\textbf{Đáp án: Máy tính bo mạch đơn.}}\par
{\footnotesize\color{sourcegray}\textit{Nguồn:} GK100 — câu 18.\par}
\vspace{3pt}\par\noindent{\color{black!18}\rule{\linewidth}{0.3pt}}\par\medskip

\needspace{5\baselineskip}
\textbf{Câu 19.} Raspberry Pi dùng bộ xử lý nào?\par
{\color{answer}\textbf{Đáp án: Bộ xử lý ARM.}}\par
{\footnotesize\color{sourcegray}\textit{Nguồn:} GK100 — câu 19.\par}
\vspace{3pt}\par\noindent{\color{black!18}\rule{\linewidth}{0.3pt}}\par\medskip

\needspace{5\baselineskip}
\textbf{Câu 20.} Raspberry Pi dùng hệ điều hành nào?\par
{\color{answer}\textbf{Đáp án: Linux.}}\par
{\footnotesize\color{sourcegray}\textit{Nguồn:} GK100 — câu 20.\par}
\vspace{3pt}\par\noindent{\color{black!18}\rule{\linewidth}{0.3pt}}\par\medskip

\needspace{5\baselineskip}
\textbf{Câu 21.} Raspberry Pi 4 Model B có mấy cổng kết nối USB\par
{\color{answer}\textbf{Đáp án: 4 cổng USB (2 × USB 3.0 và 2 × USB 2.0).}}\par
{\footnotesize\color{sourcegray}\textit{Nguồn:} GK100 — câu 21.\par}
{\footnotesize\color{sourcegray}\textit{Ghi chú AI:} Đề gốc ghi 3; sửa theo thông số chính thức của Raspberry Pi 4 Model B: \url{https://www.raspberrypi.com/products/raspberry-pi-4-model-b/specifications/}\par}
\vspace{3pt}\par\noindent{\color{black!18}\rule{\linewidth}{0.3pt}}\par\medskip

\needspace{5\baselineskip}
\textbf{Câu 22.} Cổng GPIO trong Raspberry Pi dùng để làm gì?\par
{\color{answer}\textbf{Đáp án: Tất cả các đáp án.}}\par
{\footnotesize\color{sourcegray}\textit{Nguồn:} GK100 — câu 22.\par}
\vspace{3pt}\par\noindent{\color{black!18}\rule{\linewidth}{0.3pt}}\par\medskip

\needspace{5\baselineskip}
\textbf{Câu 23.} Trong Raspberry Pi, lệnh nào để khởi động lại hệ thống?\par
{\color{answer}\textbf{Đáp án: reboot.}}\par
{\footnotesize\color{sourcegray}\textit{Nguồn:} GK100 — câu 23.\par}
\vspace{3pt}\par\noindent{\color{black!18}\rule{\linewidth}{0.3pt}}\par\medskip

\needspace{5\baselineskip}
\textbf{Câu 24.} Ưu điểm của bảng cắm dây không mối hàn là gì?\par
{\color{answer}\textbf{Đáp án: Không cần mối hàn và dễ dàng thay đổi mạch điện.}}\par
{\footnotesize\color{sourcegray}\textit{Nguồn:} GK100 — câu 24.\par}
\vspace{3pt}\par\noindent{\color{black!18}\rule{\linewidth}{0.3pt}}\par\medskip

\needspace{5\baselineskip}
\textbf{Câu 25.} Mạch điện là gì?\par
{\color{answer}\textbf{Đáp án: Đoạn đường mà dòng điện chạy qua được.}}\par
{\footnotesize\color{sourcegray}\textit{Nguồn:} GK100 — câu 25.\par}
\vspace{3pt}\par\noindent{\color{black!18}\rule{\linewidth}{0.3pt}}\par\medskip

\needspace{5\baselineskip}
\textbf{Câu 26.} Một loại bảng mẫu tạo thành đường đi cho dòng điện, được đặt dưới một khung nhựa có các lỗ cách nhau 2,54 mm (0,1 inch) và không cần hàn là gì?\par
{\color{answer}\textbf{Đáp án: Breadboard (Bảng cắm dây).}}\par
{\footnotesize\color{sourcegray}\textit{Nguồn:} GK100 — câu 26; gộp bản trùng OT50:C27.\par}
{\footnotesize\color{sourcegray}\textit{Ghi chú AI:} Nguồn trùng có cách ghi đáp án khác: OT50 C27 \ensuremath{\rightarrow} Breadboard\par}
\vspace{3pt}\par\noindent{\color{black!18}\rule{\linewidth}{0.3pt}}\par\medskip

\needspace{5\baselineskip}
\textbf{Câu 27.} Trạng thái High-Z trong điện tử có nghĩa là gì?\par
{\color{answer}\textbf{Đáp án: Trạng thái đầu ra có trở kháng rất cao.}}\par
{\footnotesize\color{sourcegray}\textit{Nguồn:} GK100 — câu 27.\par}
\vspace{3pt}\par\noindent{\color{black!18}\rule{\linewidth}{0.3pt}}\par\medskip

\needspace{5\baselineskip}
\textbf{Câu 28.} Cảm biến siêu âm để làm gì?\par
{\color{answer}\textbf{Đáp án: Đo khoảng cách.}}\par
{\footnotesize\color{sourcegray}\textit{Nguồn:} GK100 — câu 28.\par}
\vspace{3pt}\par\noindent{\color{black!18}\rule{\linewidth}{0.3pt}}\par\medskip

\needspace{5\baselineskip}
\textbf{Câu 29.} GPIO có điện áp đầu ra bao nhiêu?\par
{\color{answer}\textbf{Đáp án: 3.3V.}}\par
{\footnotesize\color{sourcegray}\textit{Nguồn:} GK100 — câu 29.\par}
\vspace{3pt}\par\noindent{\color{black!18}\rule{\linewidth}{0.3pt}}\par\medskip

\needspace{5\baselineskip}
\textbf{Câu 30.} Trong LED 7 đoạn, chân MOSI để làm gì?\par
{\color{answer}\textbf{Đáp án: Truyền dữ liệu nối tiếp từ bộ điều khiển SPI (master) tới IC ngoại vi (slave).}}\par
{\footnotesize\color{sourcegray}\textit{Nguồn:} GK100 — câu 30.\par}
{\footnotesize\color{sourcegray}\textit{Ghi chú AI:} Đề gốc ghi 'Nhận dữ liệu'; chuẩn hóa theo ý nghĩa MOSI (Master Out, Slave In). Chiều nhìn từ ngoại vi là nhận dữ liệu; vì vậy đề gốc chưa nêu rõ góc nhìn.\par}
\vspace{3pt}\par\noindent{\color{black!18}\rule{\linewidth}{0.3pt}}\par\medskip

\needspace{5\baselineskip}
\textbf{Câu 31.} Trong LED 7 đoạn, để hiển thị số 1 thì thanh nào sáng?\par
{\color{answer}\textbf{Đáp án: Các đoạn b và c của LED 7 đoạn.}}\par
{\footnotesize\color{sourcegray}\textit{Nguồn:} GK100 — câu 31.\par}
{\footnotesize\color{sourcegray}\textit{Ghi chú AI:} Viết rõ tên hai thanh sáng khi hiển thị số 1; đề gốc ghi 'BC'.\par}
\vspace{3pt}\par\noindent{\color{black!18}\rule{\linewidth}{0.3pt}}\par\medskip

\needspace{5\baselineskip}
\textbf{Câu 32.} Cảm biến độ ẩm sử dụng loại vật liệu nào?\par
{\color{answer}\textbf{Đáp án: Vật liệu tổng hợp.}}\par
{\footnotesize\color{sourcegray}\textit{Nguồn:} GK100 — câu 32.\par}
\vspace{3pt}\par\noindent{\color{black!18}\rule{\linewidth}{0.3pt}}\par\medskip

\needspace{5\baselineskip}
\textbf{Câu 33.} Khái niệm tần số là gì?\par
{\color{answer}\textbf{Đáp án: Số lần sóng chạy không tắt dao động trong một giây.}}\par
{\footnotesize\color{sourcegray}\textit{Nguồn:} GK100 — câu 33.\par}
\vspace{3pt}\par\noindent{\color{black!18}\rule{\linewidth}{0.3pt}}\par\medskip

\needspace{5\baselineskip}
\textbf{Câu 34.} Bộ khuyếch đại thuật toán có chức năng là gì?\par
{\color{answer}\textbf{Đáp án: Khuếch đại tuyến tính tín hiệu tương tự.}}\par
{\footnotesize\color{sourcegray}\textit{Nguồn:} GK100 — câu 34.\par}
\vspace{3pt}\par\noindent{\color{black!18}\rule{\linewidth}{0.3pt}}\par\medskip

\needspace{5\baselineskip}
\textbf{Câu 35.} Đây là tụ điện gì?\par
{\color{answer}\textbf{Đáp án: Tụ điện gốm (theo khóa đáp án PDF; cần đối chiếu ảnh gốc).}}\par
{\footnotesize\color{sourcegray}\textit{Nguồn:} GK100 — câu 35.\par}
{\footnotesize\color{sourcegray}\textit{Ghi chú AI:} Câu hỏi phụ thuộc hình minh họa, bản Markdown không tái tạo hình; chưa thẩm định độc lập.\par}
\vspace{3pt}\par\noindent{\color{black!18}\rule{\linewidth}{0.3pt}}\par\medskip

\needspace{5\baselineskip}
\textbf{Câu 36.} Điốt có chức năng gì trong mạch điện tử?\par
{\color{answer}\textbf{Đáp án: Chỉnh lưu và chỉnh biến.}}\par
{\footnotesize\color{sourcegray}\textit{Nguồn:} GK100 — câu 36.\par}
\vspace{3pt}\par\noindent{\color{black!18}\rule{\linewidth}{0.3pt}}\par\medskip

\needspace{5\baselineskip}
\textbf{Câu 37.} Loại điốt nào được sử dụng để chuyển mạch tốc độ cao như trong bộ trộn vi sóng và có đặc điểm là độ sụt áp thuận rất thấp ở mức 0,2V?\par
{\color{answer}\textbf{Đáp án: Điốt rào cản Schottky.}}\par
{\footnotesize\color{sourcegray}\textit{Nguồn:} GK100 — câu 37.\par}
\vspace{3pt}\par\noindent{\color{black!18}\rule{\linewidth}{0.3pt}}\par\medskip

\needspace{5\baselineskip}
\textbf{Câu 38.} Ký hiệu sau của thiết bị gì?\par
{\color{answer}\textbf{Đáp án: Khuếch đại thuật toán (theo khóa đáp án PDF; cần đối chiếu ảnh gốc).}}\par
{\footnotesize\color{sourcegray}\textit{Nguồn:} GK100 — câu 38.\par}
{\footnotesize\color{sourcegray}\textit{Ghi chú AI:} Câu hỏi phụ thuộc ký hiệu trong hình, bản Markdown không tái tạo hình; chưa thẩm định độc lập.\par}
\vspace{3pt}\par\noindent{\color{black!18}\rule{\linewidth}{0.3pt}}\par\medskip

\needspace{5\baselineskip}
\textbf{Câu 39.} Đơn vị đo điện áp trong một sơ đồ mạch điện là gì?\par
{\color{answer}\textbf{Đáp án: Vôn (V).}}\par
{\footnotesize\color{sourcegray}\textit{Nguồn:} GK100 — câu 39.\par}
\vspace{3pt}\par\noindent{\color{black!18}\rule{\linewidth}{0.3pt}}\par\medskip

\needspace{5\baselineskip}
\textbf{Câu 40.} Đơn vị đo dòng điện trong một sơ đồ mạch điện là gì?\par
{\color{answer}\textbf{Đáp án: Ampe (A).}}\par
{\footnotesize\color{sourcegray}\textit{Nguồn:} GK100 — câu 40.\par}
\vspace{3pt}\par\noindent{\color{black!18}\rule{\linewidth}{0.3pt}}\par\medskip

\needspace{5\baselineskip}
\textbf{Câu 41.} Trong một sơ đồ mạch điện, ký hiệu nào đại diện cho điện trở?\par
{\color{answer}\textbf{Đáp án: Một đường zigzag.}}\par
{\footnotesize\color{sourcegray}\textit{Nguồn:} GK100 — câu 41.\par}
\vspace{3pt}\par\noindent{\color{black!18}\rule{\linewidth}{0.3pt}}\par\medskip

\needspace{5\baselineskip}
\textbf{Câu 42.} Trạng thái nào sau đây không phải là trạng thái của cổng đệm ba trạng thái?\par
{\color{answer}\textbf{Đáp án: Trạng thái trung gian (Intermediate).}}\par
{\footnotesize\color{sourcegray}\textit{Nguồn:} GK100 — câu 42.\par}
\vspace{3pt}\par\noindent{\color{black!18}\rule{\linewidth}{0.3pt}}\par\medskip

\needspace{5\baselineskip}
\textbf{Câu 43.} Ai là người phát triển và duy trì giao diện đơn giản cho thiết bị GPIO với Raspberry Pi?\par
{\color{answer}\textbf{Đáp án: Ben Nuttall và Dave Jones.}}\par
{\footnotesize\color{sourcegray}\textit{Nguồn:} GK100 — câu 43.\par}
\vspace{3pt}\par\noindent{\color{black!18}\rule{\linewidth}{0.3pt}}\par\medskip

\needspace{5\baselineskip}
\textbf{Câu 44.} Một trong những lý do chính khiến Raspberry Pi trở nên phổ biến là gì?\par
{\color{answer}\textbf{Đáp án: 40-pin GPIO header.}}\par
{\footnotesize\color{sourcegray}\textit{Nguồn:} GK100 — câu 44.\par}
\vspace{3pt}\par\noindent{\color{black!18}\rule{\linewidth}{0.3pt}}\par\medskip

\needspace{5\baselineskip}
\textbf{Câu 45.} Trong HC-SR04, chân Trig để làm gì?\par
{\color{answer}\textbf{Đáp án: Chân TRIG nhận xung kích hoạt để HC-SR04 bắt đầu phát sóng siêu âm và đo khoảng cách.}}\par
{\footnotesize\color{sourcegray}\textit{Nguồn:} GK100 — câu 45.\par}
{\footnotesize\color{sourcegray}\textit{Ghi chú AI:} Đề gốc ghi 'Truyền siêu âm', diễn đạt không chính xác chức năng chân TRIG.\par}
\vspace{3pt}\par\noindent{\color{black!18}\rule{\linewidth}{0.3pt}}\par\medskip

\needspace{5\baselineskip}
\textbf{Câu 46.} Câu lệnh time = random(2,5) dùng để làm gì?\par
{\color{answer}\textbf{Đáp án: Tạo giá trị ngẫu nhiên trong khoảng từ 2 đến 5.}}\par
{\footnotesize\color{sourcegray}\textit{Nguồn:} GK100 — câu 46.\par}
\vspace{3pt}\par\noindent{\color{black!18}\rule{\linewidth}{0.3pt}}\par\medskip

\section*{III. Git, Linux, lập trình và công cụ}
\addcontentsline{toc}{section}{III. Git, Linux, lập trình và công cụ}

\needspace{5\baselineskip}
\textbf{Câu 47.} Quản lý phiên bản là gì?\par
{\color{answer}\textbf{Đáp án: Một hệ thống để theo dõi các phiên bản của một tài liệu.}}\par
{\footnotesize\color{sourcegray}\textit{Nguồn:} GK100 — câu 47.\par}
\vspace{3pt}\par\noindent{\color{black!18}\rule{\linewidth}{0.3pt}}\par\medskip

\needspace{5\baselineskip}
\textbf{Câu 48.} Git là gì?\par
{\color{answer}\textbf{Đáp án: Một hệ thống quản lý phiên bản phân tán.}}\par
{\footnotesize\color{sourcegray}\textit{Nguồn:} GK100 — câu 48.\par}
\vspace{3pt}\par\noindent{\color{black!18}\rule{\linewidth}{0.3pt}}\par\medskip

\needspace{5\baselineskip}
\textbf{Câu 49.} Lệnh nào trong Git được sử dụng để tạo một kho lưu trữ mới?\par
{\color{answer}\textbf{Đáp án: git init.}}\par
{\footnotesize\color{sourcegray}\textit{Nguồn:} GK100 — câu 49.\par}
\vspace{3pt}\par\noindent{\color{black!18}\rule{\linewidth}{0.3pt}}\par\medskip

\needspace{5\baselineskip}
\textbf{Câu 50.} Lệnh nào trong Git được sử dụng để đưa thay đổi vào kho lưu trữ?\par
{\color{answer}\textbf{Đáp án: git commit.}}\par
{\footnotesize\color{sourcegray}\textit{Nguồn:} GK100 — câu 50.\par}
\vspace{3pt}\par\noindent{\color{black!18}\rule{\linewidth}{0.3pt}}\par\medskip

\needspace{5\baselineskip}
\textbf{Câu 51.} Lệnh nào trong Git được sử dụng để cập nhật thay đổi từ kho lưu trữ từ xa?\par
{\color{answer}\textbf{Đáp án: git pull.}}\par
{\footnotesize\color{sourcegray}\textit{Nguồn:} GK100 — câu 51.\par}
\vspace{3pt}\par\noindent{\color{black!18}\rule{\linewidth}{0.3pt}}\par\medskip

\needspace{5\baselineskip}
\textbf{Câu 52.} Lệnh nào trong Git được sử dụng để kết hợp các thay đổi từ một nhánh khác?\par
{\color{answer}\textbf{Đáp án: git merge.}}\par
{\footnotesize\color{sourcegray}\textit{Nguồn:} GK100 — câu 52.\par}
\vspace{3pt}\par\noindent{\color{black!18}\rule{\linewidth}{0.3pt}}\par\medskip

\needspace{5\baselineskip}
\textbf{Câu 53.} GitHub là gì?\par
{\color{answer}\textbf{Đáp án: Một nền tảng lưu trữ mã nguồn trực tuyến sử dụng Git.}}\par
{\footnotesize\color{sourcegray}\textit{Nguồn:} GK100 — câu 53.\par}
\vspace{3pt}\par\noindent{\color{black!18}\rule{\linewidth}{0.3pt}}\par\medskip

\needspace{5\baselineskip}
\textbf{Câu 54.} Lệnh nào trong Git được sử dụng để xem lịch sử commit của một kho lưu trữ?\par
{\color{answer}\textbf{Đáp án: git log.}}\par
{\footnotesize\color{sourcegray}\textit{Nguồn:} GK100 — câu 54.\par}
\vspace{3pt}\par\noindent{\color{black!18}\rule{\linewidth}{0.3pt}}\par\medskip

\needspace{5\baselineskip}
\textbf{Câu 55.} Lệnh nào trong Git được sử dụng để sao chép một kho lưu trữ từ xa về máy tính của bạn?\par
{\color{answer}\textbf{Đáp án: git clone.}}\par
{\footnotesize\color{sourcegray}\textit{Nguồn:} GK100 — câu 55.\par}
\vspace{3pt}\par\noindent{\color{black!18}\rule{\linewidth}{0.3pt}}\par\medskip

\needspace{5\baselineskip}
\textbf{Câu 56.} Tại sao việc sử dụng hệ thống quản lý phiên bản lại quan trọng trong phát triển phần mềm?\par
{\color{answer}\textbf{Đáp án: Giúp theo dõi và quản lý các thay đổi trong mã nguồn.}}\par
{\footnotesize\color{sourcegray}\textit{Nguồn:} GK100 — câu 56.\par}
\vspace{3pt}\par\noindent{\color{black!18}\rule{\linewidth}{0.3pt}}\par\medskip

\needspace{5\baselineskip}
\textbf{Câu 57.} LED là gì?\par
{\color{answer}\textbf{Đáp án: Linh kiện bán dẫn phát quang.}}\par
{\footnotesize\color{sourcegray}\textit{Nguồn:} GK100 — câu 57.\par}
\vspace{3pt}\par\noindent{\color{black!18}\rule{\linewidth}{0.3pt}}\par\medskip

\needspace{5\baselineskip}
\textbf{Câu 58.} Ưu điểm của đèn LED so với bóng sợi đốt không bao gồm:\par
{\color{answer}\textbf{Đáp án: Khả năng phát sáng không giới hạn.}}\par
{\footnotesize\color{sourcegray}\textit{Nguồn:} GK100 — câu 58.\par}
\vspace{3pt}\par\noindent{\color{black!18}\rule{\linewidth}{0.3pt}}\par\medskip

\needspace{5\baselineskip}
\textbf{Câu 59.} Nhược điểm của đèn LED bao gồm?\par
{\color{answer}\textbf{Đáp án: Nguồn điện thường là 1 chiều - DC.}}\par
{\footnotesize\color{sourcegray}\textit{Nguồn:} GK100 — câu 59.\par}
\vspace{3pt}\par\noindent{\color{black!18}\rule{\linewidth}{0.3pt}}\par\medskip

\needspace{5\baselineskip}
\textbf{Câu 60.} Dòng điện chạy qua LED theo chiều?\par
{\color{answer}\textbf{Đáp án: Từ Anode (cực dương) \ensuremath{\rightarrow} Cathode (cực âm).}}\par
{\footnotesize\color{sourcegray}\textit{Nguồn:} GK100 — câu 60.\par}
\vspace{3pt}\par\noindent{\color{black!18}\rule{\linewidth}{0.3pt}}\par\medskip

\needspace{5\baselineskip}
\textbf{Câu 61.} Trên Raspberry Pi, thư viện nào được sử dụng để kết nối và điều khiển LED?\par
{\color{answer}\textbf{Đáp án: GPIO Zero.}}\par
{\footnotesize\color{sourcegray}\textit{Nguồn:} GK100 — câu 61.\par}
\vspace{3pt}\par\noindent{\color{black!18}\rule{\linewidth}{0.3pt}}\par\medskip

\needspace{5\baselineskip}
\textbf{Câu 62.} Để chạy chương trình Python led1.py thì dùng lệnh nào?\par
{\color{answer}\textbf{Đáp án: python3 led1.py.}}\par
{\footnotesize\color{sourcegray}\textit{Nguồn:} GK100 — câu 62.\par}
\vspace{3pt}\par\noindent{\color{black!18}\rule{\linewidth}{0.3pt}}\par\medskip

\needspace{5\baselineskip}
\textbf{Câu 63.} Câu nào dưới đây mở một file để đọc trong Python?\par
{\color{answer}\textbf{Đáp án: open('file.txt', 'r').}}\par
{\footnotesize\color{sourcegray}\textit{Nguồn:} GK100 — câu 63.\par}
\vspace{3pt}\par\noindent{\color{black!18}\rule{\linewidth}{0.3pt}}\par\medskip

\needspace{5\baselineskip}
\textbf{Câu 64.} Phương thức nào dưới đây dùng để đọc toàn bộ nội dung của một file?\par
{\color{answer}\textbf{Đáp án: file.read().}}\par
{\footnotesize\color{sourcegray}\textit{Nguồn:} GK100 — câu 64.\par}
\vspace{3pt}\par\noindent{\color{black!18}\rule{\linewidth}{0.3pt}}\par\medskip

\needspace{5\baselineskip}
\textbf{Câu 65.} Điều gì xảy ra nếu bạn mở một file bằng chế độ 'w'?\par
{\color{answer}\textbf{Đáp án: File sẽ được mở để ghi và dữ liệu cũ trong file sẽ bị xóa..}}\par
{\footnotesize\color{sourcegray}\textit{Nguồn:} GK100 — câu 65.\par}
\vspace{3pt}\par\noindent{\color{black!18}\rule{\linewidth}{0.3pt}}\par\medskip

\needspace{5\baselineskip}
\textbf{Câu 66.} Hệ điều hành (Operating System) là gì?\par
{\color{answer}\textbf{Đáp án: Phần mềm hệ thống.}}\par
{\footnotesize\color{sourcegray}\textit{Nguồn:} GK100 — câu 66.\par}
\vspace{3pt}\par\noindent{\color{black!18}\rule{\linewidth}{0.3pt}}\par\medskip

\needspace{5\baselineskip}
\textbf{Câu 67.} Chức năng chính của hệ điều hành là gì?\par
{\color{answer}\textbf{Đáp án: Quản lý tài nguyên hệ thống.}}\par
{\footnotesize\color{sourcegray}\textit{Nguồn:} GK100 — câu 67.\par}
\vspace{3pt}\par\noindent{\color{black!18}\rule{\linewidth}{0.3pt}}\par\medskip

\needspace{5\baselineskip}
\textbf{Câu 68.} Phần quan trọng nhất của hệ điều hành Linux là?\par
{\color{answer}\textbf{Đáp án: Kernel.}}\par
{\footnotesize\color{sourcegray}\textit{Nguồn:} GK100 — câu 68.\par}
\vspace{3pt}\par\noindent{\color{black!18}\rule{\linewidth}{0.3pt}}\par\medskip

\needspace{5\baselineskip}
\textbf{Câu 69.} Hệ điều hành đa nhiệm là gì?\par
{\color{answer}\textbf{Đáp án: Hệ điều hành chạy nhiều ứng dụng cùng lúc.}}\par
{\footnotesize\color{sourcegray}\textit{Nguồn:} GK100 — câu 69.\par}
\vspace{3pt}\par\noindent{\color{black!18}\rule{\linewidth}{0.3pt}}\par\medskip

\needspace{5\baselineskip}
\textbf{Câu 70.} Shell nào được coi là shell mặc định phổ biến nhất trên các hệ thống Linux hiện nay?\par
{\color{answer}\textbf{Đáp án: Bash (Bourne Again Shell).}}\par
{\footnotesize\color{sourcegray}\textit{Nguồn:} GK100 — câu 70.\par}
\vspace{3pt}\par\noindent{\color{black!18}\rule{\linewidth}{0.3pt}}\par\medskip

\needspace{5\baselineskip}
\textbf{Câu 71.} Chương trình soạn thảo văn bản nào là phổ biến nhất trong hệ điều hành Linux?\par
{\color{answer}\textbf{Đáp án: Vim (theo đáp án đề gốc).}}\par
{\footnotesize\color{sourcegray}\textit{Nguồn:} GK100 — câu 71.\par}
{\footnotesize\color{sourcegray}\textit{Ghi chú AI:} Khẳng định 'phổ biến nhất' tùy cộng đồng và thời điểm; không có đáp án khách quan duy nhất.\par}
\vspace{3pt}\par\noindent{\color{black!18}\rule{\linewidth}{0.3pt}}\par\medskip

\needspace{5\baselineskip}
\textbf{Câu 72.} Trong hệ thống Linux, user nào có quyền cao nhất?\par
{\color{answer}\textbf{Đáp án: root.}}\par
{\footnotesize\color{sourcegray}\textit{Nguồn:} GK100 — câu 72.\par}
\vspace{3pt}\par\noindent{\color{black!18}\rule{\linewidth}{0.3pt}}\par\medskip

\needspace{5\baselineskip}
\textbf{Câu 73.} Shell trong hệ điều hành có chức năng chính là gì?\par
{\color{answer}\textbf{Đáp án: Tiếp nhận và thực thi các lệnh từ người dùng..}}\par
{\footnotesize\color{sourcegray}\textit{Nguồn:} GK100 — câu 73.\par}
\vspace{3pt}\par\noindent{\color{black!18}\rule{\linewidth}{0.3pt}}\par\medskip

\needspace{5\baselineskip}
\textbf{Câu 74.} Lệnh echo được sử dụng để làm gì?\par
{\color{answer}\textbf{Đáp án: Để hiển thị một thông báo hoặc chuỗi văn bản trên màn hình..}}\par
{\footnotesize\color{sourcegray}\textit{Nguồn:} GK100 — câu 74.\par}
\vspace{3pt}\par\noindent{\color{black!18}\rule{\linewidth}{0.3pt}}\par\medskip

\needspace{5\baselineskip}
\textbf{Câu 75.} Lệnh nào dùng để kết hợp đầu ra của lệnh cat file1 và cat file2 vào một tập tin output.txt?\par
{\color{answer}\textbf{Đáp án: cat file1 file2 \textgreater{} output.txt.}}\par
{\footnotesize\color{sourcegray}\textit{Nguồn:} GK100 — câu 75.\par}
\vspace{3pt}\par\noindent{\color{black!18}\rule{\linewidth}{0.3pt}}\par\medskip

\needspace{5\baselineskip}
\textbf{Câu 76.} Biến môi trường HOME trong shell đại diện cho gì?\par
{\color{answer}\textbf{Đáp án: Đường dẫn đến thư mục chính của người dùng hiện tại..}}\par
{\footnotesize\color{sourcegray}\textit{Nguồn:} GK100 — câu 76.\par}
\vspace{3pt}\par\noindent{\color{black!18}\rule{\linewidth}{0.3pt}}\par\medskip

\needspace{5\baselineskip}
\textbf{Câu 77.} Để hủy bỏ lệnh đang thi hành, ta thực hiện?\par
{\color{answer}\textbf{Đáp án: Ctrl + C để gửi tín hiệu ngắt (SIGINT) cho lệnh chạy ở foreground.}}\par
{\footnotesize\color{sourcegray}\textit{Nguồn:} GK100 — câu 77.\par}
{\footnotesize\color{sourcegray}\textit{Ghi chú AI:} Đề gốc ghi Ctrl + Z, nhưng Ctrl + Z gửi SIGTSTP để tạm dừng tiến trình, không phải hủy/ngắt.\par}
\vspace{3pt}\par\noindent{\color{black!18}\rule{\linewidth}{0.3pt}}\par\medskip

\needspace{5\baselineskip}
\textbf{Câu 78.} IDE là viết tắt của từ gì?\par
{\color{answer}\textbf{Đáp án: Integrated Development Environment.}}\par
{\footnotesize\color{sourcegray}\textit{Nguồn:} GK100 — câu 78.\par}
\vspace{3pt}\par\noindent{\color{black!18}\rule{\linewidth}{0.3pt}}\par\medskip

\needspace{5\baselineskip}
\textbf{Câu 79.} Tại sao các lập trình viên thường sử dụng IDE?\par
{\color{answer}\textbf{Đáp án: Tất cả các phương án trên.}}\par
{\footnotesize\color{sourcegray}\textit{Nguồn:} GK100 — câu 79.\par}
\vspace{3pt}\par\noindent{\color{black!18}\rule{\linewidth}{0.3pt}}\par\medskip

\needspace{5\baselineskip}
\textbf{Câu 80.} Trong Nano, để thoát khỏi trình soạn thảo văn bản, bạn sử dụng phím tắt nào?\par
{\color{answer}\textbf{Đáp án: Ctrl + X.}}\par
{\footnotesize\color{sourcegray}\textit{Nguồn:} GK100 — câu 80.\par}
\vspace{3pt}\par\noindent{\color{black!18}\rule{\linewidth}{0.3pt}}\par\medskip

\needspace{5\baselineskip}
\textbf{Câu 81.} Trong Vim, để lưu và thoát khỏi trình soạn thảo, bạn sử dụng phím tắt nào?\par
{\color{answer}\textbf{Đáp án: :wq.}}\par
{\footnotesize\color{sourcegray}\textit{Nguồn:} GK100 — câu 81.\par}
\vspace{3pt}\par\noindent{\color{black!18}\rule{\linewidth}{0.3pt}}\par\medskip

\needspace{5\baselineskip}
\textbf{Câu 82.} Để tạo một thư mục mới, bạn sử dụng lệnh nào?\par
{\color{answer}\textbf{Đáp án: mkdir.}}\par
{\footnotesize\color{sourcegray}\textit{Nguồn:} GK100 — câu 82.\par}
\vspace{3pt}\par\noindent{\color{black!18}\rule{\linewidth}{0.3pt}}\par\medskip

\needspace{5\baselineskip}
\textbf{Câu 83.} Lệnh nào dùng để hiển thị các tiến trình đang chạy của người dùng?\par
{\color{answer}\textbf{Đáp án: ps.}}\par
{\footnotesize\color{sourcegray}\textit{Nguồn:} GK100 — câu 83.\par}
\vspace{3pt}\par\noindent{\color{black!18}\rule{\linewidth}{0.3pt}}\par\medskip

\needspace{5\baselineskip}
\textbf{Câu 84.} Lệnh nào dùng để đổi tên hoặc di chuyển tệp và thư mục?\par
{\color{answer}\textbf{Đáp án: mv.}}\par
{\footnotesize\color{sourcegray}\textit{Nguồn:} GK100 — câu 84.\par}
\vspace{3pt}\par\noindent{\color{black!18}\rule{\linewidth}{0.3pt}}\par\medskip

\needspace{5\baselineskip}
\textbf{Câu 85.} Lệnh nào hiển thị đường dẫn của thư mục hiện tại?\par
{\color{answer}\textbf{Đáp án: pwd.}}\par
{\footnotesize\color{sourcegray}\textit{Nguồn:} GK100 — câu 85.\par}
\vspace{3pt}\par\noindent{\color{black!18}\rule{\linewidth}{0.3pt}}\par\medskip

\section*{IV. Mạng máy tính, HTTP và SQL}
\addcontentsline{toc}{section}{IV. Mạng máy tính, HTTP và SQL}

\needspace{5\baselineskip}
\textbf{Câu 86.} Mạng máy tính là?\par
{\color{answer}\textbf{Đáp án: Các máy tính kết nối qua môi trường truyền tin và trao đổi thông tin theo một kiến trúc mạng xác định.}}\par
{\footnotesize\color{sourcegray}\textit{Nguồn:} GK100 — câu 86.\par}
\vspace{3pt}\par\noindent{\color{black!18}\rule{\linewidth}{0.3pt}}\par\medskip

\needspace{5\baselineskip}
\textbf{Câu 87.} Website hoạt động dựa trên mô hình gì?\par
{\color{answer}\textbf{Đáp án: Client-Server.}}\par
{\footnotesize\color{sourcegray}\textit{Nguồn:} GK100 — câu 87.\par}
\vspace{3pt}\par\noindent{\color{black!18}\rule{\linewidth}{0.3pt}}\par\medskip

\needspace{5\baselineskip}
\textbf{Câu 88.} Để gửi mail chúng ta sử dụng giao thức gì?\par
{\color{answer}\textbf{Đáp án: SMTP.}}\par
{\footnotesize\color{sourcegray}\textit{Nguồn:} GK100 — câu 88.\par}
\vspace{3pt}\par\noindent{\color{black!18}\rule{\linewidth}{0.3pt}}\par\medskip

\needspace{5\baselineskip}
\textbf{Câu 89.} HTTP viết tắt bởi?\par
{\color{answer}\textbf{Đáp án: HyperText Transfer Protocol.}}\par
{\footnotesize\color{sourcegray}\textit{Nguồn:} GK100 — câu 89.\par}
\vspace{3pt}\par\noindent{\color{black!18}\rule{\linewidth}{0.3pt}}\par\medskip

\needspace{5\baselineskip}
\textbf{Câu 90.} Dung lượng mạng (Bandwidth) được đo bằng đơn vị nào?\par
{\color{answer}\textbf{Đáp án: Bits per second (bps).}}\par
{\footnotesize\color{sourcegray}\textit{Nguồn:} GK100 — câu 90.\par}
\vspace{3pt}\par\noindent{\color{black!18}\rule{\linewidth}{0.3pt}}\par\medskip

\needspace{5\baselineskip}
\textbf{Câu 91.} Loại cáp nào có tốc độ truyền dữ liệu cao nhất?\par
{\color{answer}\textbf{Đáp án: Cáp quang (Fiber optic cable).}}\par
{\footnotesize\color{sourcegray}\textit{Nguồn:} GK100 — câu 91.\par}
\vspace{3pt}\par\noindent{\color{black!18}\rule{\linewidth}{0.3pt}}\par\medskip

\needspace{5\baselineskip}
\textbf{Câu 92.} Mã trạng thái 200 nghĩa là gì?\par
{\color{answer}\textbf{Đáp án: Yêu cầu của khách hàng đã được máy chủ xử lý thành công..}}\par
{\footnotesize\color{sourcegray}\textit{Nguồn:} GK100 — câu 92.\par}
\vspace{3pt}\par\noindent{\color{black!18}\rule{\linewidth}{0.3pt}}\par\medskip

\needspace{5\baselineskip}
\textbf{Câu 93.} Mã trạng thái 503 nghĩa là gì?\par
{\color{answer}\textbf{Đáp án: Service Unavailable.}}\par
{\footnotesize\color{sourcegray}\textit{Nguồn:} GK100 — câu 93.\par}
\vspace{3pt}\par\noindent{\color{black!18}\rule{\linewidth}{0.3pt}}\par\medskip

\needspace{5\baselineskip}
\textbf{Câu 94.} Wi-Fi sử dụng các tiêu chuẩn nào sau đây để truyền dữ liệu không dây?\par
{\color{answer}\textbf{Đáp án: IEEE 802.11.}}\par
{\footnotesize\color{sourcegray}\textit{Nguồn:} GK100 — câu 94.\par}
\vspace{3pt}\par\noindent{\color{black!18}\rule{\linewidth}{0.3pt}}\par\medskip

\needspace{5\baselineskip}
\textbf{Câu 95.} Dữ liệu được gửi qua phương thức GET sẽ được đính kèm ở đâu trong yêu cầu HTTP?\par
{\color{answer}\textbf{Đáp án: Trong URL của yêu cầu.}}\par
{\footnotesize\color{sourcegray}\textit{Nguồn:} GK100 — câu 95.\par}
\vspace{3pt}\par\noindent{\color{black!18}\rule{\linewidth}{0.3pt}}\par\medskip

\needspace{5\baselineskip}
\textbf{Câu 96.} Khi nào nên sử dụng phương thức GET thay vì POST?\par
{\color{answer}\textbf{Đáp án: Khi yêu cầu có thể được lưu vào bộ nhớ đệm hoặc được chia sẻ qua URL.}}\par
{\footnotesize\color{sourcegray}\textit{Nguồn:} GK100 — câu 96.\par}
\vspace{3pt}\par\noindent{\color{black!18}\rule{\linewidth}{0.3pt}}\par\medskip

\needspace{5\baselineskip}
\textbf{Câu 97.} Ưu điểm chính của giao thức truyền thông nối tiếp so với truyền thông song song là gì?\par
{\color{answer}\textbf{Đáp án: Chi phí thấp hơn..}}\par
{\footnotesize\color{sourcegray}\textit{Nguồn:} GK100 — câu 97.\par}
\vspace{3pt}\par\noindent{\color{black!18}\rule{\linewidth}{0.3pt}}\par\medskip

\needspace{5\baselineskip}
\textbf{Câu 98.} Lệnh nào được sử dụng để tính tổng một cột trong SQL?\par
{\color{answer}\textbf{Đáp án: SUM.}}\par
{\footnotesize\color{sourcegray}\textit{Nguồn:} GK100 — câu 98.\par}
\vspace{3pt}\par\noindent{\color{black!18}\rule{\linewidth}{0.3pt}}\par\medskip

\needspace{5\baselineskip}
\textbf{Câu 99.} SQL viết tắt của gì?\par
{\color{answer}\textbf{Đáp án: Structured Query Language.}}\par
{\footnotesize\color{sourcegray}\textit{Nguồn:} GK100 — câu 99.\par}
\vspace{3pt}\par\noindent{\color{black!18}\rule{\linewidth}{0.3pt}}\par\medskip

\needspace{5\baselineskip}
\textbf{Câu 100.} Lệnh nào được sử dụng để cập nhật dữ liệu trong một bảng trong SQL?\par
{\color{answer}\textbf{Đáp án: UPDATE.}}\par
{\footnotesize\color{sourcegray}\textit{Nguồn:} GK100 — câu 100.\par}
\vspace{3pt}\par\noindent{\color{black!18}\rule{\linewidth}{0.3pt}}\par\medskip

\section*{V. Tổng quan, kiến trúc và ứng dụng IoT}
\addcontentsline{toc}{section}{V. Tổng quan, kiến trúc và ứng dụng IoT}

\needspace{5\baselineskip}
\textbf{Câu 101.} Số lượng thiết bị IoT do IOT Analytics dự báo đến năm 2030 đạt được là?\par
{\color{answer}\textbf{Đáp án: 39 tỷ thiết bị IoT được kết nối vào năm 2030 (theo dự báo IoT Analytics công bố 28/10/2025).}}\par
{\footnotesize\color{sourcegray}\textit{Nguồn:} OT50 — câu 1.\par}
{\footnotesize\color{sourcegray}\textit{Ghi chú AI:} Đề gốc ghi 29 tỷ; đây là dự báo cũ. Sửa theo https://iot-analytics.com/Number-Connected-Iot-Devices/\par}
\vspace{3pt}\par\noindent{\color{black!18}\rule{\linewidth}{0.3pt}}\par\medskip

\needspace{5\baselineskip}
\textbf{Câu 102.} Thuật ngữ "Internet of Things" do ai sáng tạo ra?\par
{\color{answer}\textbf{Đáp án: Kevin Ashton.}}\par
{\footnotesize\color{sourcegray}\textit{Nguồn:} OT50 — câu 2.\par}
\vspace{3pt}\par\noindent{\color{black!18}\rule{\linewidth}{0.3pt}}\par\medskip

\needspace{5\baselineskip}
\textbf{Câu 103.} Trong mô hình IoT 4 lớp, lớp nào chịu trách nhiệm trực quan hóa dữ liệu?\par
{\color{answer}\textbf{Đáp án: Lớp giao diện (theo bộ lựa chọn của đề).}}\par
{\footnotesize\color{sourcegray}\textit{Nguồn:} OT50 — câu 3.\par}
{\footnotesize\color{sourcegray}\textit{Ghi chú AI:} Lưu ý: slide Chương 1 mô tả lớp ứng dụng cung cấp giao diện; tên 'lớp giao diện' không phải một lớp độc lập trong kiến trúc IoT bốn lớp đang dùng.\par}
\vspace{3pt}\par\noindent{\color{black!18}\rule{\linewidth}{0.3pt}}\par\medskip

\needspace{5\baselineskip}
\textbf{Câu 104.} Chức năng chính của lớp nhận thức (lớp thiết bị/lớp vật lý) là?\par
{\color{answer}\textbf{Đáp án: Thu thập dữ liệu từ môi trường và tác động môi trường.}}\par
{\footnotesize\color{sourcegray}\textit{Nguồn:} OT50 — câu 4.\par}
\vspace{3pt}\par\noindent{\color{black!18}\rule{\linewidth}{0.3pt}}\par\medskip

\needspace{5\baselineskip}
\textbf{Câu 105.} Vấn đề nào không phải là thách thức của IoT?\par
{\color{answer}\textbf{Đáp án: Mức tiêu thụ điện năng thấp.}}\par
{\footnotesize\color{sourcegray}\textit{Nguồn:} OT50 — câu 5.\par}
\vspace{3pt}\par\noindent{\color{black!18}\rule{\linewidth}{0.3pt}}\par\medskip

\needspace{5\baselineskip}
\textbf{Câu 106.} Lớp cảm nhận trong mô hình IoT bao gồm?\par
{\color{answer}\textbf{Đáp án: Tất cả đều đúng.}}\par
{\footnotesize\color{sourcegray}\textit{Nguồn:} OT50 — câu 6.\par}
\vspace{3pt}\par\noindent{\color{black!18}\rule{\linewidth}{0.3pt}}\par\medskip

\needspace{5\baselineskip}
\textbf{Câu 107.} Máy vi tính ra đời năm nào?\par
{\color{answer}\textbf{Đáp án: 1946.}}\par
{\footnotesize\color{sourcegray}\textit{Nguồn:} OT50 — câu 7.\par}
\vspace{3pt}\par\noindent{\color{black!18}\rule{\linewidth}{0.3pt}}\par\medskip

\needspace{5\baselineskip}
\textbf{Câu 108.} Chức năng chính của lớp xử lý là?\par
{\color{answer}\textbf{Đáp án: Lưu trữ và xử lý dữ liệu.}}\par
{\footnotesize\color{sourcegray}\textit{Nguồn:} OT50 — câu 8.\par}
\vspace{3pt}\par\noindent{\color{black!18}\rule{\linewidth}{0.3pt}}\par\medskip

\needspace{5\baselineskip}
\textbf{Câu 109.} Khái niệm nào liên quan đến IoT?\par
{\color{answer}\textbf{Đáp án: Tất cả đều đúng.}}\par
{\footnotesize\color{sourcegray}\textit{Nguồn:} OT50 — câu 9.\par}
\vspace{3pt}\par\noindent{\color{black!18}\rule{\linewidth}{0.3pt}}\par\medskip

\needspace{5\baselineskip}
\textbf{Câu 110.} Mô hình IoT 4 lớp không bao gồm lớp nào?\par
{\color{answer}\textbf{Đáp án: Lớp giao diện.}}\par
{\footnotesize\color{sourcegray}\textit{Nguồn:} OT50 — câu 10.\par}
\vspace{3pt}\par\noindent{\color{black!18}\rule{\linewidth}{0.3pt}}\par\medskip

\needspace{5\baselineskip}
\textbf{Câu 111.} Kiến trúc IoT gồm 4 thành phần cơ bản?\par
{\color{answer}\textbf{Đáp án: Things + Gateway (Network) + Cloud + Services.}}\par
{\footnotesize\color{sourcegray}\textit{Nguồn:} OT50 — câu 11.\par}
{\footnotesize\color{sourcegray}\textit{Ghi chú AI:} Đề gốc ghi Things + Gateway + User Interface + Services, trái với slide Chapter 1, mục IoT Architecture (4 basic components).\par}
\vspace{3pt}\par\noindent{\color{black!18}\rule{\linewidth}{0.3pt}}\par\medskip

\needspace{5\baselineskip}
\textbf{Câu 112.} Điện thoại ra đời năm nào?\par
{\color{answer}\textbf{Đáp án: 1876.}}\par
{\footnotesize\color{sourcegray}\textit{Nguồn:} OT50 — câu 12.\par}
\vspace{3pt}\par\noindent{\color{black!18}\rule{\linewidth}{0.3pt}}\par\medskip

\needspace{5\baselineskip}
\textbf{Câu 113.} Một trong những công nghệ nền tảng của Công nghiệp 4.0 là?\par
{\color{answer}\textbf{Đáp án: IIoT.}}\par
{\footnotesize\color{sourcegray}\textit{Nguồn:} OT50 — câu 13.\par}
\vspace{3pt}\par\noindent{\color{black!18}\rule{\linewidth}{0.3pt}}\par\medskip

\needspace{5\baselineskip}
\textbf{Câu 114.} Thuật ngữ Internet of Things lần đầu xuất hiện năm nào?\par
{\color{answer}\textbf{Đáp án: 1999.}}\par
{\footnotesize\color{sourcegray}\textit{Nguồn:} OT50 — câu 14.\par}
\vspace{3pt}\par\noindent{\color{black!18}\rule{\linewidth}{0.3pt}}\par\medskip

\needspace{5\baselineskip}
\textbf{Câu 115.} IoT là gì?\par
{\color{answer}\textbf{Đáp án: Tất cả đều đúng.}}\par
{\footnotesize\color{sourcegray}\textit{Nguồn:} OT50 — câu 15.\par}
\vspace{3pt}\par\noindent{\color{black!18}\rule{\linewidth}{0.3pt}}\par\medskip

\needspace{5\baselineskip}
\textbf{Câu 116.} Chức năng chính của lớp ứng dụng?\par
{\color{answer}\textbf{Đáp án: Cung cấp dịch vụ cho ứng dụng.}}\par
{\footnotesize\color{sourcegray}\textit{Nguồn:} OT50 — câu 16.\par}
\vspace{3pt}\par\noindent{\color{black!18}\rule{\linewidth}{0.3pt}}\par\medskip

\needspace{5\baselineskip}
\textbf{Câu 117.} Học phần cung cấp kiến thức gì?\par
{\color{answer}\textbf{Đáp án: Tất cả đều đúng.}}\par
{\footnotesize\color{sourcegray}\textit{Nguồn:} OT50 — câu 17.\par}
\vspace{3pt}\par\noindent{\color{black!18}\rule{\linewidth}{0.3pt}}\par\medskip

\needspace{5\baselineskip}
\textbf{Câu 118.} Internet hiện đang trong giai đoạn?\par
{\color{answer}\textbf{Đáp án: Internet of Things (theo sơ đồ bốn giai đoạn trong Chapter 1).}}\par
{\footnotesize\color{sourcegray}\textit{Nguồn:} OT50 — câu 18.\par}
{\footnotesize\color{sourcegray}\textit{Ghi chú AI:} Mô hình bốn giai đoạn là khung trình bày trong giáo trình, không phải khẳng định phân kỳ lịch sử có hiệu lực vĩnh viễn.\par}
\vspace{3pt}\par\noindent{\color{black!18}\rule{\linewidth}{0.3pt}}\par\medskip

\needspace{5\baselineskip}
\textbf{Câu 119.} Internet ra đời năm nào?\par
{\color{answer}\textbf{Đáp án: 1969.}}\par
{\footnotesize\color{sourcegray}\textit{Nguồn:} OT50 — câu 19.\par}
\vspace{3pt}\par\noindent{\color{black!18}\rule{\linewidth}{0.3pt}}\par\medskip

\needspace{5\baselineskip}
\textbf{Câu 120.} WWW và HTML ra đời năm nào?\par
{\color{answer}\textbf{Đáp án: 1990.}}\par
{\footnotesize\color{sourcegray}\textit{Nguồn:} OT50 — câu 20.\par}
\vspace{3pt}\par\noindent{\color{black!18}\rule{\linewidth}{0.3pt}}\par\medskip

\needspace{5\baselineskip}
\textbf{Câu 121.} Công nghệ nền tảng dẫn đến hình thành IoT?\par
{\color{answer}\textbf{Đáp án: RFID.}}\par
{\footnotesize\color{sourcegray}\textit{Nguồn:} OT50 — câu 21.\par}
\vspace{3pt}\par\noindent{\color{black!18}\rule{\linewidth}{0.3pt}}\par\medskip

\needspace{5\baselineskip}
\textbf{Câu 122.} Mô hình IoT 5 lớp bổ sung lớp?\par
{\color{answer}\textbf{Đáp án: Kinh doanh.}}\par
{\footnotesize\color{sourcegray}\textit{Nguồn:} OT50 — câu 22.\par}
\vspace{3pt}\par\noindent{\color{black!18}\rule{\linewidth}{0.3pt}}\par\medskip

\needspace{5\baselineskip}
\textbf{Câu 123.} Thành phần thu thập dữ liệu môi trường?\par
{\color{answer}\textbf{Đáp án: Thiết bị cảm biến.}}\par
{\footnotesize\color{sourcegray}\textit{Nguồn:} OT50 — câu 23.\par}
\vspace{3pt}\par\noindent{\color{black!18}\rule{\linewidth}{0.3pt}}\par\medskip

\section*{VI. Arduino, cảm biến và truyền thông IoT}
\addcontentsline{toc}{section}{VI. Arduino, cảm biến và truyền thông IoT}

\needspace{5\baselineskip}
\textbf{Câu 124.} Toán tử khác nhau trong Arduino?\par
{\color{answer}\textbf{Đáp án: !=.}}\par
{\footnotesize\color{sourcegray}\textit{Nguồn:} OT50 — câu 24.\par}
\vspace{3pt}\par\noindent{\color{black!18}\rule{\linewidth}{0.3pt}}\par\medskip

\needspace{5\baselineskip}
\textbf{Câu 125.} Hàm căn bậc hai trong Arduino?\par
{\color{answer}\textbf{Đáp án: sqrt().}}\par
{\footnotesize\color{sourcegray}\textit{Nguồn:} OT50 — câu 25.\par}
\vspace{3pt}\par\noindent{\color{black!18}\rule{\linewidth}{0.3pt}}\par\medskip

\needspace{5\baselineskip}
\textbf{Câu 126.} analogRead(pin) dùng để?\par
{\color{answer}\textbf{Đáp án: Đọc chân tương tự.}}\par
{\footnotesize\color{sourcegray}\textit{Nguồn:} OT50 — câu 26.\par}
\vspace{3pt}\par\noindent{\color{black!18}\rule{\linewidth}{0.3pt}}\par\medskip

\needspace{5\baselineskip}
\textbf{Câu 127.} Hàm chỉ chạy một lần?\par
{\color{answer}\textbf{Đáp án: setup.}}\par
{\footnotesize\color{sourcegray}\textit{Nguồn:} OT50 — câu 28.\par}
\vspace{3pt}\par\noindent{\color{black!18}\rule{\linewidth}{0.3pt}}\par\medskip

\needspace{5\baselineskip}
\textbf{Câu 128.} Không phải cấu trúc điều khiển?\par
{\color{answer}\textbf{Đáp án: setup.}}\par
{\footnotesize\color{sourcegray}\textit{Nguồn:} OT50 — câu 29.\par}
\vspace{3pt}\par\noindent{\color{black!18}\rule{\linewidth}{0.3pt}}\par\medskip

\needspace{5\baselineskip}
\textbf{Câu 129.} Cảm biến tích cực?\par
{\color{answer}\textbf{Đáp án: Cảm biến siêu âm.}}\par
{\footnotesize\color{sourcegray}\textit{Nguồn:} OT50 — câu 30.\par}
\vspace{3pt}\par\noindent{\color{black!18}\rule{\linewidth}{0.3pt}}\par\medskip

\needspace{5\baselineskip}
\textbf{Câu 130.} Cảm biến nào sau đây là cảm biến thụ động (passive sensor)?\par
{\color{answer}\textbf{Đáp án: Cảm biến áp suất.}}\par
{\footnotesize\color{sourcegray}\textit{Nguồn:} OT50 — câu 31.\par}
\vspace{3pt}\par\noindent{\color{black!18}\rule{\linewidth}{0.3pt}}\par\medskip

\needspace{5\baselineskip}
\textbf{Câu 131.} Thiết bị nào trong các thiết bị sau là cơ cấu chấp hành?\par
{\color{answer}\textbf{Đáp án: LCD.}}\par
{\footnotesize\color{sourcegray}\textit{Nguồn:} OT50 — câu 32.\par}
\vspace{3pt}\par\noindent{\color{black!18}\rule{\linewidth}{0.3pt}}\par\medskip

\needspace{5\baselineskip}
\textbf{Câu 132.} Sản phẩm nào sau đây là một vi máy tính?\par
{\color{answer}\textbf{Đáp án: Raspberry Pi.}}\par
{\footnotesize\color{sourcegray}\textit{Nguồn:} OT50 — câu 33.\par}
\vspace{3pt}\par\noindent{\color{black!18}\rule{\linewidth}{0.3pt}}\par\medskip

\needspace{5\baselineskip}
\textbf{Câu 133.} Dòng lệnh pinMode(pin, mode) dùng để làm gì?\par
{\color{answer}\textbf{Đáp án: Thiết lập chế độ ra (OUTPUT) hoặc vào (INPUT) cho chân được chỉ định.}}\par
{\footnotesize\color{sourcegray}\textit{Nguồn:} OT50 — câu 34.\par}
\vspace{3pt}\par\noindent{\color{black!18}\rule{\linewidth}{0.3pt}}\par\medskip

\needspace{5\baselineskip}
\textbf{Câu 134.} Thành phần nào trong hệ thống IoT thực hiện chức năng chuyển đổi các tín hiệu điện thành chuyển động cơ học hoặc các biến đổi vật lý khác (như áp suất, nhiệt độ...)?\par
{\color{answer}\textbf{Đáp án: Thiết bị truyền động.}}\par
{\footnotesize\color{sourcegray}\textit{Nguồn:} OT50 — câu 35.\par}
\vspace{3pt}\par\noindent{\color{black!18}\rule{\linewidth}{0.3pt}}\par\medskip

\needspace{5\baselineskip}
\textbf{Câu 135.} Giao thức nào sau đây thuộc lớp ứng dụng trong mô hình truyền thông?\par
{\color{answer}\textbf{Đáp án: HTTP.}}\par
{\footnotesize\color{sourcegray}\textit{Nguồn:} OT50 — câu 36.\par}
\vspace{3pt}\par\noindent{\color{black!18}\rule{\linewidth}{0.3pt}}\par\medskip

\needspace{5\baselineskip}
\textbf{Câu 136.} Giao thức nào sau đây hoạt động theo mô hình Request/Response?\par
{\color{answer}\textbf{Đáp án: HTTP.}}\par
{\footnotesize\color{sourcegray}\textit{Nguồn:} OT50 — câu 37.\par}
\vspace{3pt}\par\noindent{\color{black!18}\rule{\linewidth}{0.3pt}}\par\medskip

\needspace{5\baselineskip}
\textbf{Câu 137.} Nhận dạng khuôn mặt là kiểu giao tiếp nào trong hệ thống IoT?\par
{\color{answer}\textbf{Đáp án: H2M (Human-to-Machine).}}\par
{\footnotesize\color{sourcegray}\textit{Nguồn:} OT50 — câu 38; gộp bản trùng OT50:C39.\par}
{\footnotesize\color{sourcegray}\textit{Ghi chú AI:} Đề gốc ghi M2H; slide Chapter 3, 'Types of Communications in IoT' liệt kê facial recognition là ví dụ H2M.\par}
\vspace{3pt}\par\noindent{\color{black!18}\rule{\linewidth}{0.3pt}}\par\medskip

\needspace{5\baselineskip}
\textbf{Câu 138.} Giao thức truyền thông nào sau đây có phạm vi hoạt động xa và khả năng truyền tải tốc độ cao?\par
{\color{answer}\textbf{Đáp án: 4G.}}\par
{\footnotesize\color{sourcegray}\textit{Nguồn:} OT50 — câu 40.\par}
\vspace{3pt}\par\noindent{\color{black!18}\rule{\linewidth}{0.3pt}}\par\medskip

\needspace{5\baselineskip}
\textbf{Câu 139.} Địa chỉ nào sau đây là địa chỉ IP?\par
{\color{answer}\textbf{Đáp án: 225.252.255.52.}}\par
{\footnotesize\color{sourcegray}\textit{Nguồn:} OT50 — câu 41.\par}
\vspace{3pt}\par\noindent{\color{black!18}\rule{\linewidth}{0.3pt}}\par\medskip

\needspace{5\baselineskip}
\textbf{Câu 140.} Các mô hình truyền thông Kết nối hệ thống mở (OSI) và TCP/IP tương ứng được chia thành bao nhiêu lớp chức năng?\par
{\color{answer}\textbf{Đáp án: 7 và 4.}}\par
{\footnotesize\color{sourcegray}\textit{Nguồn:} OT50 — câu 42.\par}
\vspace{3pt}\par\noindent{\color{black!18}\rule{\linewidth}{0.3pt}}\par\medskip

\needspace{5\baselineskip}
\textbf{Câu 141.} Giao thức truyền thông là gì?\par
{\color{answer}\textbf{Đáp án: Một bộ quy tắc cho phép hai hoặc nhiều thiết bị thiết lập liên kết và trao đổi dữ liệu.}}\par
{\footnotesize\color{sourcegray}\textit{Nguồn:} OT50 — câu 43.\par}
\vspace{3pt}\par\noindent{\color{black!18}\rule{\linewidth}{0.3pt}}\par\medskip

\needspace{5\baselineskip}
\textbf{Câu 142.} Trong mô hình truyền thông IoT, lớp nào thực hiện chức năng gán địa chỉ và định tuyến cho các gói tin?\par
{\color{answer}\textbf{Đáp án: Lớp mạng.}}\par
{\footnotesize\color{sourcegray}\textit{Nguồn:} OT50 — câu 44.\par}
\vspace{3pt}\par\noindent{\color{black!18}\rule{\linewidth}{0.3pt}}\par\medskip

\needspace{5\baselineskip}
\textbf{Câu 143.} Giao thức truyền thông nào sau đây có kiểu hình mạng Mesh?\par
{\color{answer}\textbf{Đáp án: Zigbee.}}\par
{\footnotesize\color{sourcegray}\textit{Nguồn:} OT50 — câu 45.\par}
\vspace{3pt}\par\noindent{\color{black!18}\rule{\linewidth}{0.3pt}}\par\medskip

\needspace{5\baselineskip}
\textbf{Câu 144.} Trong mô hình truyền thông của hệ thống IoT, giao thức WiFi thuộc lớp nào?\par
{\color{answer}\textbf{Đáp án: Lớp vật lý và MAC.}}\par
{\footnotesize\color{sourcegray}\textit{Nguồn:} OT50 — câu 46.\par}
\vspace{3pt}\par\noindent{\color{black!18}\rule{\linewidth}{0.3pt}}\par\medskip

\needspace{5\baselineskip}
\textbf{Câu 145.} Hình thức giao tiếp mà không có sự tương tác với con người là?\par
{\color{answer}\textbf{Đáp án: M2M.}}\par
{\footnotesize\color{sourcegray}\textit{Nguồn:} OT50 — câu 47.\par}
\vspace{3pt}\par\noindent{\color{black!18}\rule{\linewidth}{0.3pt}}\par\medskip

\needspace{5\baselineskip}
\textbf{Câu 146.} Giao thức truyền thông nào sau đây chỉ có phạm vi hoạt động gần và khả năng truyền tải tốc độ thấp?\par
{\color{answer}\textbf{Đáp án: Zigbee (theo cách phân loại trong đề gốc).}}\par
{\footnotesize\color{sourcegray}\textit{Nguồn:} OT50 — câu 48.\par}
{\footnotesize\color{sourcegray}\textit{Ghi chú AI:} Bluetooth/BLE cũng là công nghệ tầm ngắn tốc độ thấp; câu gốc không đủ điều kiện để phân biệt duy nhất.\par}
\vspace{3pt}\par\noindent{\color{black!18}\rule{\linewidth}{0.3pt}}\par\medskip

\needspace{5\baselineskip}
\textbf{Câu 147.} Giao thức lớp ứng dụng nào sau đây hoạt động theo mô hình Publish/Subscribe?\par
{\color{answer}\textbf{Đáp án: MQTT.}}\par
{\footnotesize\color{sourcegray}\textit{Nguồn:} OT50 — câu 49.\par}
\vspace{3pt}\par\noindent{\color{black!18}\rule{\linewidth}{0.3pt}}\par\medskip

\needspace{5\baselineskip}
\textbf{Câu 148.} Giao thức truyền thông nào sau đây chỉ có phạm vi hoạt động rất xa nhưng khả năng truyền tải tốc độ thấp?\par
{\color{answer}\textbf{Đáp án: LoRaWAN.}}\par
{\footnotesize\color{sourcegray}\textit{Nguồn:} OT50 — câu 50.\par}
\vspace{3pt}\par\noindent{\color{black!18}\rule{\linewidth}{0.3pt}}\par\medskip

\section*{VII. Câu hỏi bổ sung sau đối chiếu lịch sử quiz}
\addcontentsline{toc}{section}{VII. Câu hỏi bổ sung sau đối chiếu lịch sử quiz}

\needspace{5\baselineskip}
\textbf{Câu 149.} Trong giai đoạn 'Internet of Peoples', Internet thường được dùng để làm gì?\par
{\color{answer}\textbf{Đáp án: Xem video, trò chuyện trực tuyến, chơi game trực tuyến và các hoạt động tương tác giữa người dùng..}}\par
{\footnotesize\color{sourcegray}\textit{Nguồn:} Chapter 1: Immersive Experiences; nhiều bản Quiz W1..\par}
{\footnotesize\color{sourcegray}\textit{Ghi chú AI:} TỰ TỔNG HỢP (AI): Các phương án riêng đều phù hợp nên trả lời tổng hợp.\par}
\vspace{3pt}\par\noindent{\color{black!18}\rule{\linewidth}{0.3pt}}\par\medskip

\needspace{5\baselineskip}
\textbf{Câu 150.} Theo mô hình được giảng dạy trong Chương 1, Internet phát triển qua bao nhiêu giai đoạn?\par
{\color{answer}\textbf{Đáp án: 4: Connectivity, Networked Economy, Immersive Experiences, Internet of Things..}}\par
{\footnotesize\color{sourcegray}\textit{Nguồn:} Chapter 1, 'The History of the Internet', trang trình chiếu số 5–6; Quiz W1..\par}
{\footnotesize\color{sourcegray}\textit{Ghi chú AI:} TỰ XÁC ĐỊNH (AI) từ slide học phần.\par}
\vspace{3pt}\par\noindent{\color{black!18}\rule{\linewidth}{0.3pt}}\par\medskip

\needspace{5\baselineskip}
\textbf{Câu 151.} Theo mô hình bốn giai đoạn trong Chương 1, mục nào KHÔNG phải tên một giai đoạn số hóa của Internet: số hóa truy cập, số hóa kinh doanh, số hóa thông tin, số hóa thế giới?\par
{\color{answer}\textbf{Đáp án: Số hóa thông tin. Các giai đoạn trong slide: digitize access, digitize business, digitize interactions, digitize the world..}}\par
{\footnotesize\color{sourcegray}\textit{Nguồn:} Chapter 1, slide số 6; nhiều Quiz W1..\par}
{\footnotesize\color{sourcegray}\textit{Ghi chú AI:} TỰ XÁC ĐỊNH (AI).\par}
\vspace{3pt}\par\noindent{\color{black!18}\rule{\linewidth}{0.3pt}}\par\medskip

\needspace{5\baselineskip}
\textbf{Câu 152.} Ứng dụng IoT nào thể hiện trực tiếp lĩnh vực thành phố thông minh trong các lựa chọn giao thông, thương mại, công nghiệp và y tế?\par
{\color{answer}\textbf{Đáp án: Giảm thiểu tắc nghẽn giao thông bằng hệ thống giao thông thông minh..}}\par
{\footnotesize\color{sourcegray}\textit{Nguồn:} Nhiều bản Quiz W1; Chapter 1 phần IoT Applications..\par}
{\footnotesize\color{sourcegray}\textit{Ghi chú AI:} TỰ XÁC ĐỊNH (AI) từ ngữ cảnh đáp án.\par}
\vspace{3pt}\par\noindent{\color{black!18}\rule{\linewidth}{0.3pt}}\par\medskip

\needspace{5\baselineskip}
\textbf{Câu 153.} Khái niệm cốt lõi của IoT theo Chương 1 là gì?\par
{\color{answer}\textbf{Đáp án: Đối tượng thông minh (Smart Objects) và môi trường thông minh (Smart Environments)..}}\par
{\footnotesize\color{sourcegray}\textit{Nguồn:} Chapter 1, 'IoT Core Concepts'; nhiều bản Quiz W1..\par}
{\footnotesize\color{sourcegray}\textit{Ghi chú AI:} TỰ XÁC ĐỊNH (AI): Các bản quiz thay đổi phương án, vì vậy không dùng đáp án 'Tất cả đều đúng' chung cho mọi phiên bản.\par}
\vspace{3pt}\par\noindent{\color{black!18}\rule{\linewidth}{0.3pt}}\par\medskip

\needspace{5\baselineskip}
\textbf{Câu 154.} Trong kiến trúc IoT bốn lớp, lớp mạng có chức năng chính gì?\par
{\color{answer}\textbf{Đáp án: Kết nối các thiết bị, các mạng và truyền tải dữ liệu giữa chúng..}}\par
{\footnotesize\color{sourcegray}\textit{Nguồn:} Chapter 1, 'Service-Oriented Based Architecture' slide số 35; nhiều bản Quiz W1..\par}
{\footnotesize\color{sourcegray}\textit{Ghi chú AI:} TỰ XÁC ĐỊNH (AI), trích ý từ giáo trình.\par}
\vspace{3pt}\par\noindent{\color{black!18}\rule{\linewidth}{0.3pt}}\par\medskip

\needspace{5\baselineskip}
\textbf{Câu 155.} Trong danh sách các thiết bị nút bấm, microphone, radar và LCD, đâu là thiết bị đầu ra thay vì cảm biến/đầu vào?\par
{\color{answer}\textbf{Đáp án: LCD (màn hình hiển thị); không phải cơ cấu chấp hành cơ học theo nghĩa hẹp..}}\par
{\footnotesize\color{sourcegray}\textit{Nguồn:} Quiz W2 (PDF dscdvsv); Chapter 2 về Sensors/Actuators..\par}
{\footnotesize\color{sourcegray}\textit{Ghi chú AI:} TỰ HIỆU ĐÍNH (AI): Câu gốc tự mâu thuẫn ('cơ cấu chấp hành' nhưng còn yêu cầu 'chọn một đáp án sai'); đã viết lại để có đáp án xác định.\par}
\vspace{3pt}\par\noindent{\color{black!18}\rule{\linewidth}{0.3pt}}\par\medskip

\needspace{5\baselineskip}
\textbf{Câu 156.} Trong các thiết bị microphone, LED, LCD và loa kỹ thuật số, thiết bị nào là cảm biến?\par
{\color{answer}\textbf{Đáp án: Microphone (chuyển dao động âm thanh thành tín hiệu điện).}}\par
{\footnotesize\color{sourcegray}\textit{Nguồn:} Quiz W2, PDF \texttt{dscdvsv(2).pdf}, câu 7; Chapter 2: IoT Hardware.\par}
{\footnotesize\color{sourcegray}\textit{Ghi chú AI:} TỰ XÁC ĐỊNH/CHUẨN HÓA (AI). Bỏ chỉ dẫn mâu thuẫn 'hãy chọn một đáp án sai' trong văn bản gốc; phân loại theo chức năng kỹ thuật của thiết bị.\par}
\vspace{3pt}\par\noindent{\color{black!18}\rule{\linewidth}{0.3pt}}\par\medskip

\needspace{5\baselineskip}
\textbf{Câu 157.} Một chương trình (đơn vị mã) được tải lên và chạy trên module Arduino được gọi là gì?\par
{\color{answer}\textbf{Đáp án: Sketch.}}\par
{\footnotesize\color{sourcegray}\textit{Nguồn:} Quiz W2 (N01, N02, N04), cụ thể Quiz W2 – lớp N04, câu 1 (1/1 điểm, lựa chọn “Sketch”); PDF \texttt{Quiz W2\_ Cơ sở và thiết kế hệ thốngSSSSS IoT-1-1-26(N04).pdf}, trang 1.\par}
{\footnotesize\color{sourcegray}\textit{Ghi chú AI:} TỰ CHUẨN HÓA (AI) phần câu hỏi: tài liệu gốc có thêm chỉ dẫn “Không cần giải thích, hãy chọn 1 đáp án sai!” mâu thuẫn với mục tiêu chọn đáp án chính xác. Đã loại bỏ chỉ dẫn đó; đáp án “Sketch” được xác nhận bởi bản quiz đạt 1/1 điểm, không phải đáp án tự đoán.\par}
\vspace{3pt}\par\noindent{\color{black!18}\rule{\linewidth}{0.3pt}}\par\medskip

\end{document}
